\documentclass[a4paper]{article}
\usepackage[normalem]{ulem}

% impostazioni generali
%Tutti gli usepackage vanno qui
\usepackage{geometry}
\usepackage[italian]{babel}
\usepackage[utf8]{inputenc}
\usepackage{tabularx}
\usepackage{longtable}
\usepackage{hyperref}
\usepackage{enumitem}
\usepackage{array} 
\usepackage{booktabs}
\newcolumntype{M}[1]{>{\centering\arraybackslash}m{#1}}
\usepackage[toc]{appendix}

\hypersetup{
	colorlinks=true,
	linkcolor=blue,
	filecolor=magenta,
	urlcolor=blue,
}
% Numerazione figure
\let\counterwithout\relax
\let\counterwithin\relax
\usepackage{chngcntr}

% distanziare elenco delle figure e delle tabelle
\usepackage{tocbasic}
\DeclareTOCStyleEntry[numwidth=3.5em]{tocline}{figure}% for figure entries
\DeclareTOCStyleEntry[numwidth=3.5em]{tocline}{table}% for table entries


\counterwithin{table}{subsection}
\counterwithin{figure}{subsection}

\usepackage[bottom]{footmisc}
\usepackage{fancyhdr}
\setcounter{secnumdepth}{4}
\usepackage{amsmath, amssymb}
\usepackage{array}
\usepackage{graphicx}

\usepackage{ifthen}

\usepackage{float}
\restylefloat{table}

\usepackage{layouts}
\usepackage{url}
\usepackage{comment}
\usepackage{eurosym}

\usepackage{lastpage}
\usepackage{layouts}
\usepackage{eurosym}

\geometry{a4paper,top=3cm,bottom=4cm,left=2.5cm,right=2.5cm}

%Comandi di impaginazione uguale per tutti i documenti
\pagestyle{fancy}
\lhead{\includegraphics[scale=0.1]{../../../template/images/logo_no_motto.jpeg}}
%Titolo del documento
\rhead{\doctitle{}}
%\rfoot{\thepage}
\cfoot{Pagina \thepage\ di \pageref{LastPage}}
\setlength{\headheight}{35pt}
\setcounter{tocdepth}{5}
\setcounter{secnumdepth}{5}
\renewcommand{\footrulewidth}{0.4pt}

% multirow per tabelle
\usepackage{multirow}

% Permette tabelle su più pagine
%\usepackage{longtable}


% colore di sfondo per le celle
\usepackage[table]{xcolor}

%COMANDI TABELLE
\newcommand{\rowcolorhead}{\rowcolor[HTML]{007c95}}
\newcommand{\cellcolorhead}{\cellcolor[HTML]{007c95}}
\newcommand{\hlinetable}{\arrayrulecolor[HTML]{007c95}\hline}

%intestazione
% check for missing commands
\newcommand{\headertitle}[1]{\textbf{\color{white}#1}} %titolo colonna
\definecolor{pari}{HTML}{b1dae3}
\definecolor{dispari}{HTML}{d7f2f7}

% comandi glossario
\newcommand{\glo}{$_{G}$}
\newcommand{\glosp}{$_{G}$ }


%label custom
\makeatletter
\newcommand{\uclabel}[2]{%
	\protected@write \@auxout {}{\string \newlabel {#1}{{#2}{\thepage}{#2}{#1}{}} }%
	\hypertarget{#1}{#2}
}
\makeatother

%riportare pezzi di codice
\definecolor{codegray}{gray}{0.9}
\newcommand{\code}[1]{\colorbox{codegray}{\texttt{#1}}}



% dati relativi alla prima pagina
% Configurazione della pagina iniziale
\newcommand{\doctitle}{\textit{Verbale} interno 01-12-2022 }
\newcommand{\docdate}{01 Dicembre 2022}
\newcommand{\rev}{1.0.0}
\newcommand{\stato}{Approvato}
\newcommand{\uso}{Interno}
\newcommand{\approv}{Elena Pandolfo}
\newcommand{\red}{Tommaso Allegretti}
\newcommand{\ver}{Gabriele Mantoan\\& Mirko Stella}
\newcommand{\dest}{\textit{Seven Clickers}
									  \\ Prof. Vardanega Tullio 
									   \\ Prof. Cardin Riccardo}
\newcommand{\describedoc}{\textit{Verbale} riguardante il meeting tenuto il giorno 01-12-2022}


 % editare questo

\makeindex

\begin{document}
\counterwithin{table}{section}

% Prima pagina
\thispagestyle{empty}
\renewcommand{\arraystretch}{1.3}

\begin{titlepage}
	\begin{center}
		
	\includegraphics[scale = 0.40]{../../../template/images/logo.jpeg}
	\\[1cm]
	\href{mailto:7clickersgroup@gmail.com}		      	
	{\large{\textit{7clickersgroup@gmail.com} } }\\[2.5cm]
	\Huge \textbf{\doctitle} \\[1cm]
	 \large
			 \begin{tabular}{r|l}
                        \textbf{Versione} & \rev{} \\
                        \textbf{Stato} & \stato{} \\
                        \textbf{Uso} & \uso{} \\                         
                        \textbf{Approvazione} & \approv{} \\                      
                        \textbf{Redazione} & \red{} \\ 
                        \textbf{Verifica} &  \ver{} \\                         
                        \textbf{Distribuzione} & \parbox[t]{5cm}{ \dest{} }
                \end{tabular} 
                \\[3.5cm]
                \large \textbf{Descrizione} \\ \describedoc{} 
     \end{center}
\end{titlepage}

% Diario delle modifiche
\section*{Registro delle modifiche}

\newcommand{\changelogTable}[1]{
	 

\renewcommand{\arraystretch}{1.5}
\rowcolors{2}{pari}{dispari}
\begin{longtable}{ 
		>{\centering}M{0.07\textwidth} 
		>{\centering}M{0.13\textwidth}
		>{\centering}M{0.20\textwidth}
		>{\centering}M{0.17\textwidth} 
		>{\centering\arraybackslash}M{0.30\textwidth} 
		 }
	\rowcolorhead
	\headertitle{Vers.} &
	\centering \headertitle{Data} &	
	\headertitle{Autore} &
	\headertitle{Ruolo} & 
	\headertitle{Descrizione} 
	\endfirsthead	
	\endhead
	
	#1

\end{longtable}
\vspace{-2em}

}



\changelogTable{
	1.0.0 & 14-12-22 & Elena Pandolfo & Responsabile &  Approvazione documento\tabularnewline
	0.1.0 & 01-12-22 & Gabriele Mantoan \\ Mirko Stella & Verificatori & Verifica documento\tabularnewline
	0.0.1 & 09-11-22 & Rino Sincic & Analista &  Redazione documento\\
} % editare questo
\pagebreak

% Indice
{
    \hypersetup{linkcolor=black}
    \tableofcontents
}
\pagebreak

% Contenuto
\section{Introduzione}
\subsection{Scopo del documento}
Lo scopo del documento è quello di stabilire le regole che ogni componente del gruppo SevenClickers
deve rispettare per mantenere un ambiente di lavoro che mira a massimizzare
l'economicità dei processi durante il ciclo di vita del prodotto ShowRoom3D.\\
Le norme verranno inserite in modo incrementale per regolamentare le attività di progetto imminenti rimandando quelle meno urgenti a quando 
se ne presenterà la necessità.\\
Inoltre tali norme potranno subire modifiche nel tempo in modo da garantire un miglioramento continuo della qualità del lavoro svolto.
Il responsabile di progetto ha il compito di comunicare l'aggiunta di una nuova norma o la modifica di una già esistente a tutti i componenti
del gruppo e di assicurarsi che siano comprese a pieno.
\subsection{Scopo del prodotto}
Il prodotto in questione nasce dalla necessità dell'azienda SanMarco Informatica di fornire una soluzione agli sprechi
derivati dall'adozione di uno ShowRoom tradizionale proponendo uno ShowRoom3D che sia 
ugualmente o ancora più immersivo.

\section{Documentazione}
Questa sezione descrive le convenzioni,gli strumenti e le modalità con cui il gruppo si impegna a stilare la documentazione interna ed esterna relativa al progetto.

\subsubsection{Convenzioni generali}
Le convenzioni di seguito riportate vengono applicate a tutti i documenti.
Esse rendono i documenti stilati omogenei tra loro contribuendo a rendere il progetto professionale.

\paragraph{Versionamento}
Il numero di versione permette di capire lo stato in cui si trova un documento.
Un documento può trovarsi nei seguenti stati:
\begin{itemize} 
    \item \textbf{Approvato}: il documento è verificato ed approvato dal Responsabile;
    \item \textbf{Verificato}: il documento risulta verificato ma non ancora visionato dal Responsabile;
    \item \textbf{In Sviluppo}: sono presenti delle modifiche che non sono state verificate.
\end{itemize}
Il numero di versione ha il formato \textbf{X.Y.Z} dove:
\begin{itemize} 
    \item \textbf{X}: indica una versione approvata dal Responsabile, la numerazione parte da 0
    e la prima versione approvata è la 1.0.0;
    \item \textbf{Y}: indica una versione verificata dal Verificatore, la numerazione inizia da 0 e si azzera ad ogni incremento di X. La prima versione
    verificata è la 0.1.0;
    \item \textbf{Z}: indica una versione in fase di modifica da parte dei redattori che ne incrementano il numero ad ogni modifica,
    la numerazione parte da 1 e si azzera ad ogni incremento di X o Y. La prima versione modificata è la 0.0.1.
\end{itemize}

\paragraph{Struttura Generale}
Ogni documento deve presentare le seguenti sezioni nell'ordine in cui vengono presentate:
\begin{itemize} 
    \item \textbf{Intestazione}:
    contiene:
    \begin{itemize} 
        \item Logo compreso di motto;
        \item Indirizzo email di gruppo; 
        \item Titolo;
        \item Tabella contenente le informazioni generali:
        \begin{itemize}
            \item Versione;
            \item Stato;
            \item Uso;
            \item Approvazione: indica il responsabile di progetto che ha approvato il documento; 
            \item Redazione: elenco dei collaboratori che hanno partecipato alla stesura del documento;
            \item Verifica: elenco dei verificatori che hanno verificato il documento;
            \item Distribuzione: elenco delle persone o organizzazioni a cui è destinato il documento.
        \end{itemize}
        \item Breve descrizione del documento .
    \end{itemize}
    \item \textbf{Registro delle modifiche}:
    tabella che identifica ogni versione del documento indicandone:
    \begin{itemize} 
        \item Versione;
        \item Data;
        \item Autore;
        \item Ruolo;
        \item Descrizione.
    \end{itemize}
    \item \textbf{Indice}:
    elenco ordinato dei titoli dei capitoli, ovvero delle varie parti di cui si compone il documento;
    \item \textbf{Contenuto}:
    varia a seconda del tipo di documento.
\end{itemize}
\paragraph{Date}
Le date devono rispettare il seguente formato: \textbf{dd-mm-yyyy}
All'interno delle tabelle il formato deve essere il seguente: \textbf{dd-mm-yy}
\paragraph{Nomi di persona}
All'interno dei documenti i nomi di persona rispetteranno l'ordine nome seguito dal cognome della persona menzionata.

\paragraph{Elenchi puntati}
Gli elenchi puntati devono rispettare le seguenti regole:
\begin{itemize} 
    \item Ogni elemento dell'elenco deve iniziare con la lettera maiuscola;
    \item Ogni elemento dell'elenco deve terminare con ";" ad eccezione dell'ultimo elemento
    che deve terminare con "."; 
    \item Dopo i due punti la frase deve iniziare con la lettera minuscola.
\end{itemize}

\paragraph{Stile del testo}
\begin{itemize} 
    \item \textbf{Grassetto}: stile utilizzato per i titoli delle sezioni e per i primi termini degli elenchi puntati;
    \item \textbf{Corsivo}: viene utilizzato per citare il nome dei documenti, ad esempio \textit{Piano di Progetto}. 
\end{itemize}

\subsection{Verifica}
La verifica viene svolta da due verificatori prima del merge con il branch documentation.
Consiste nell'esaminare i file prodotti da chi ne ha fatto la stesura e segnalarne la non validità o 
la presenza di errori nei concetti esposti.\\
Un verificatore dovrà verificare il documento a partire dalle modifiche fatte dopo l'ultima versione verificata.
Le modifiche da verificare quindi possono essere dedotte dal registro dei cambiamenti presente in ogni documento.
Una volta controllato il documento, il primo verificatore segnalerà eventuali errori e
successivamente dovrà spuntare come approvata la Pull Request nella sezione dedicata su GitHub\textsubscript{g}. \\
A questo punto se un secondo verificatore noterà la necessità di qualche altro cambiamento da apportare, chi dovrà apportare le modifiche farà una pull in locale per
allineare il proprio branch con quello in remoto e continuare con il proprio lavoro.\\
Dopo il push delle modifiche se i file risultano corretti anche dal secondo verificatore esso aggiungerà i nomi dei verificatori all'intestazione modificando il file titlepage\_input.tex
 e creerà la riga nel registro delle modifiche, nel file changelog\_input.tex,  inserendo la nuova versione secondo le regole di versionamento e scrivendo nella colonna Autore sia il suo nome,
che quello del primo verificatore, e in descrizione "Verifica".
Dopo aver eseguito un commit sul ramo da integrare approva la Pull Request e conferma il merge secondo le norme di progetto descritte nella sezione dedicata.
\subsection{Approvazione}
L’approvazione viene svolta dal Responsabile di Progetto.\\
L'approvazione consiste nell'aprire una Pull Request di approvazione. Nel caso in cui il Responsabile di progetto riscontrasse ulteriori problematiche, segnalerà ai Verificatori le eventuali modifiche da apportare. I verificatori apporteranno tali modifiche prima di chiudere la Pull Request di approvazione e quindi di integrare i cambiamenti nel branch documentation.\\
Per i verbali si effettua questa pratica non appena il file prodotto viene verificato e quindi passa tutte le eventuali reviews create dai verificatori mentre per tutti gli altri documenti prima di una consegna.\\
Ad esito positivo di approvazione, il Responsabile di Progetto creerà una nuova riga e compilerà i campi nelle colonne corrispondenti del registro delle modifiche inserendo: l'ultima versione secondo le norme di versionamento, la data, il proprio nome, il suo ruolo e la voce "Approvazione" nell'ultima colonna.
\subsubsection{Strumenti per la stesura}
\begin{itemize} 
    \item LaTeX: è un linguaggio di marcatura per la preparazione di testi, basato sul 				  		  programma di composizione tipografica TEX.\\
	Nel branch documentation  si possono trovare i file .pdf prodotti e la cartella “latex”. La cartella latex contiene tre cartelle interne:
	\begin{itemize}
	\item esterni e interni, contengono file .tex di documentazione esterna ed interna come ad esempio i verbali o altra documentazione esterna/interna:
	\begin{itemize}
		\item la cartella config, contiene i file .tex con le parti fisse dei documenti (intestazione,registro delle modifiche,tracciamento dei temi affrontati) che vengono modificati con i dati del documento specifico 
		\item la cartella res/sections, contiene i file .tex con il contenuto vero e proprio (sezioni del documento) che viene redatto in maniera libera dal redattore
		\item un file col nome del documento pdf con estensione .tex che viene compilato per produrre il file pdf
	\end{itemize}
	 \item template, contiene file .tex di base utilizzati secondo necessità per comporre i documenti:
	 \begin{itemize}
	 \item changelox.tex è il file di template che serve per scrivere il registro delle modifiche
	 \item package.tex è il file che contiene tutti gli usepackage\textsubscript{g}
	 \item titlepage.tex è il file di template che contiene la configurazione della pagina iniziale di ogni documento  
	 \item tracking.tex è il file di template che contiene il tracciamento dei temi affrontati nel documento
	 \end{itemize}
	\end{itemize}
\end{itemize}

\section{Strumenti collaborativi}
\subsection{GitHub\textsubscript{g}}
Servizio di hosting per progetti software che implementa uno strumento di controllo versione distribuito Git\textsubscript{g}.
Oltre alla copia in remoto del repository\textsubscript{g} di progetto ogni componente del gruppo ha una propria copia in locale.\\
Per ottenere una copia del repository\textsubscript{g} ogni componente ha scaricato lo strumento Git\textsubscript{g} ed eseguendo il 
comando 'git clone' da git\textsubscript{g} bash viene creata una cartella collegata alla repository\textsubscript{g} di progetto.\\
Non sono state imposte modalità specifiche sull'interazione con il repository\textsubscript{g} remoto in modo da non sconvolgere le abitudini di lavoro di 
ciascun componente.\\
I componenti del gruppo abituati ad interagire con GitHub\textsubscript{g} da interfaccia grafica possono continuare a farne uso.
\subsubsection{Repository\textsubscript{g}}
Il repository\textsubscript{g} si può trovare all'indirizzo \textbf{\textit{https://github.com/7clickers/ShowRoom3D}} ed è pubblico. 
I collaboratori sono i componenti del gruppo SevenClickers che utilizzano il proprio account GitHub\textsubscript{g} personale per collaborare al progetto.
\subsubsection{Branching}
\textbf{Branches protetti}:
    \begin{itemize} 
        \item main: contiene le versioni di release\textsubscript{g} del software
        \item documentation: contiene i template latex\textsubscript{g} e rispettivi pdf della documentazione
    \end{itemize}
\textit{documentation}: 
I documenti presenti in documentation sono stati approvati dal Responsabile di progetto o almeno verificati dai verificatori.\\
Per integrare delle modifiche da un branch protetto ad uno libero si utilizza un branch d’appoggio creato in locale partendo dall'ultimo commit di documentation e facendone il merge con il branch che necessita delle integrazioni. In seguito il branch di appoggio verrà eliminato.\\
\textbf{Branches liberi}:
Vengono utilizzati per creare nuove funzionalità e gli sviluppatori possono effettuare i commit senza l'approvazione degli altri componenti del gruppo
in quanto ciascun componente sviluppa su un solo branch alla volta salvo casi eccezionali.
Un branch\textsubscript{g} libero avrà il nome del documento che si sviluppa su quel branch\textsubscript{g}, oppure della feature\textsubscript{g} che va ad implementare.\\
Non appena i/il file nel branch sono stati verificati ed il merge è stato fatto, il branch libero verrà eliminato.
I branch di approvazione saranno chiamati con la sintassi appr\_nomefile1\_nomefile2\_nomefile3.... a seconda dei file che verranno approvati durante il ciclo di vita del branch.
\subsubsection{Commits}
È preferibile che ogni commit abbia una singola responsabilità per cambiamento.\\
I commits non possono essere effettuati direttamente sui branch protetti ma per integrare delle aggiunte o modifiche sarà necessario aprire una Pull Request.
All'approvazione di una Pull Request tutti i commit relativi al merge verranno raggruppati in un unico commit che rispetti la struttura sintattica descritta in seguito.

I commit dovranno essere accompagnati da una descrizione solo se ritenuta indispensabile alla comprensione del commit stesso.

I messaggi di commit sui \textbf{\uppercase{branch protetti}} dovranno seguire la seguente struttura sintattica:\\

\textbf{$<$label$>$$<$\#n\_issue$>$$<$testo$>$}\\\\
dove:\\\\
\textbf{label}: può assumere i seguenti valori
\begin{itemize}
\item feat: indica che è stata implementata una nuova funzionalità
\item fix: indica che è stato risolto un bug
\item update: indica che è stata apportata una modifica che non sia fix o feat
\item test: qualsiasi cosa legata ai test
\item docs: qualsiasi cosa legata alla documentazione
\end{itemize}
\textbf{n\_issue}: indica il numero della issue a cui fa riferimento il commit (se non fa riferimento a nessuna issue viene omesso).\\\\
\textbf{testo}:indica con quale branch è stato effettuato il merge e deve rispettare la forma: merge from $<$nome branch da integrare$>$ to $<$nome branch corrente$>$ \\\\
\textbf{descrizione}: se aggiunta ad un commit deve rispondere alle domande WHAT?,WHY?,HOW? ovvero
cosa è cambiato,perchè sono stati fatti i cambiamenti,in che modo sono stati fatti i cambiamenti.\\

I messaggi di commit sui \textbf{\uppercase{branch liberi}} dovranno seguire la struttura sintattica dei branch protetti ad eccezione del testo.
Il testo dei commit sui branch liberi non è soggetto a restrizioni particolari a patto che indichi in maniera intuitiva i cambiamenti fatti
in modo che possano essere compresi anche dagli altri collaboratori.



\subsubsection{Pull Requests}
Per effettuare un merge su un branch protetto si deve aprire da GitHub\textsubscript{g} una Pull Request.
La Pull Request permette di verificare il lavoro svolto prima di integrarlo con il branch desiderato.\\
Alla creazione di una Pull Request bisogna associare:
\begin{itemize}
\item i verificatori in carica hanno il compito di trovare eventuali errori o mancanze e fornire un feedback riguardante il contenuto direttamente su GitHub\textsubscript{g} richiedendo
una review con un review comment sulla parte specifica da revisionare o con un commento generico.\\
Non sarà possibile effettuare il merge finchè tutti i commenti di revisione non saranno stati risolti e la Pull Request approvata da due verificatori
\item l’issue associata nell’opzione “Development” che verrà chiusa alla risoluzione della Pull Request
\item la Projects Board di cui fa parte
\item gli assegnatari che hanno il compito di apportare le modifiche necessarie in fase di verifica
\item le labels associate
\end{itemize}
Per i commit relativi alle Pull Requests seguire le regole descritte nella sottosezione Commits per i branch protetti.

\subsubsection{Milestone\textsubscript{g}}
Una milestone indica un traguardo intermedio significativo per il progetto.
Ad essa possono venire assegnate delle issues per verificarne il raggiungimento.
Ogni milestone ha una scadenza che viene discussa e fissata da tutto il gruppo.
Oltre alle 2 milestone + 1 milestone falcoltativa fissate dal committente il gruppo ne creerá ulteriori per scandire piú nel dettaglio i passi che ci porteranno 
ad ottenere i risultati prefissati.
Una delle milestone create dal gruppo durante il completamento della tecnology baseline relativa alla prima milestone imposta dal committente (Requirements and Tecnology Baseline)
riguarda l'implementazione di un PoC (Proof of concept).

\subsubsection{Projects Board\textsubscript{g}}
Vengono utilizzate due project board per tracciare le issues della repo.
Una project board principale utilizzata da tutti i membri del gruppo e una project board per le approvazioni utilizzata solo dal responsabile di progetto per approvare i file che richiedono l'approvazione prima di una consegna.

\begin{itemize}
	\item La projectboard\textsubscript{g} principale è suddivisa in queste sezioni:
	\begin{itemize}
		\item \textbf{Todo}: Issue\textsubscript{g} che non sono ancora state iniziate o che non sono ancora state assegnate
		\item \textbf{In Progress}: Issue\textsubscript{g} che sono state assegnate e a cui almeno un membro a cui è stata assegnata ha iniziato a lavorarci
		\item \textbf{Pull Request}: Issue\textsubscript{g} che è in fase di integrazione e necessita della verifica dei verificatori. Corrisponde all'inizio di una pull request
		\item \textbf{Done}: Issue\textsubscript{g} che sono state chiuse e che sono state verificate (se necessitano di verifica\textsubscript{g})
		\item \textbf{Approved}: Issue\textsubscript{g} con label "Da Approvare" che hanno ottenuto l'approvazione\textsubscript{g} del responsabile subito dopo la verifica.
		Per tutte le issues che richiedono un'approvazione solo al momento della consegna è stata creata la project board dedicata alle approvazioni.
	\end{itemize}
	Inoltre nella project board principale vengono registrate delle issue che non richiedono verifica, approvazione o neanche integrazione, con lo scopo di monitorare meglio il lavoro di ogni membro del team.\\
Queste issue verranno chiuse e archiviate manualmente una volta che avranno terminato la loro utilità, un esempio può essere la seguente issue:\\
\textbf{diario di bordo 21-11-22}; questa issue non necessità verifica, approvazione o integrazione perchè non è di interesse caricare il file nella repo, però è utile tracciare lo svolgimento della issue.

	\item La projectboard\textsubscript{g} riservata alle approvazioni è suddivisa in queste sezioni:
	\begin{itemize}
		\item \textbf{Todo}: Issue\textsubscript{g} che non sono ancora state iniziate dal responsabile di progetto
		\item \textbf{In Progress}: Issue\textsubscript{g}  che sono state prese in carico per essere approvate dal responsabile di progetto
		\item \textbf{Pull Request}: Issue\textsubscript{g} che è in fase di integrazione e necessita della verifica dei verificatori. Corrisponde all'inizio di una pull request
		in questo caso i verificatori dovranno solo controllare che l'intestazione e il registro delle modifiche siano stati compilati correttamente dal responsabile di progetto in quanto tutto il resto del contenuto è già stato verificato in precedenza
		\item \textbf{Approved}: Issue\textsubscript{g} che sono state approvate dal responsabile di progetto
	\end{itemize}
\end{itemize}
Le issue che si trovano nella project board approvazioni sono state create come continuazione ad issues che hanno in precedenza attraversato il work flow della project board principale fino allo stato "Done".
In questo modo il Responsabile di progetto dovrà approvare solamente file che si trovano nei branches protetti.
Vengono ricreate in questa project board per  mantenerne la tracciabilità e permetterne una più facile reperibilità futura per l'approvazione.
Questo permette anche di pulire la project board principale da tutte quelle issues che rimarrebbero inattive per molto tempo.
\\\\
Esempio di work flow issues:
\begin{enumerate}
	\item Viene creata l'issue ed inserita all'interno della sezione "To Do" della project board principale.
	\item Quando viene presa in carico l'issue viene spostata nella sezione "In Progress" fino al suo completamento.
	\item L'incaricato alla risoluzione dell'issue apre una pull request a cui assegna l'issue in modo che  venga chiusa in automatico al momento dell'integrazione con il branch di destinazione.
	\item Vengono fatte le verifiche opportune dai verificatori e le conseguenti modifiche prima di accettare la pull request.
	\item Dopo l'accettazione la pull request viene postata insieme alle issues associate nella colonna "Done".
	\item A questo punto si possono presentare due casi: 
	\begin{enumerate}
		\item I file relativi alla pull request richiedono approvazione immediata.
		Viene archiviata la pull request (e le issue associate) e creata una issue a cui viene aggiunta la label "Da approvare" con il nome del file da approvare.
		Il responsabile creerá un branch per approvare le issues con approvazione immediata aprirá una pull request per integrare nel branch di destinazione l'approvazione.
		\item I file relativi alla pull request richiedono l'approvazione prima di una consegna.
		Il responsabile archivia la pull request e ne crea una issue con il nome dei file da approvare nella project board dedicata alle approvazioni.
		NB:Se nella project board Approvazioni é giá presente un'issue di approvazione per il file da approvare non viene creata una nuova issue di approvazione ma si tiene la precedente.
		L'issue segue il work flow della project board approvazioni.
	\end{enumerate}
\end{enumerate}
\subsubsection{Issue Tracking System\textsubscript{g}}
Il gruppo utilizza l'issue tracking system\textsubscript{g} di Github\textsubscript{g} per tenere traccia delle issue\textsubscript{g}. 
Le issues\textsubscript{g} verranno determinate dal responsabile, ma la loro assegnazione verrà effettuata dai membri del gruppo, in base alla priorità delle issues\textsubscript{g}, i loro ruoli e alle loro disponibilità temporali.\\
Per marchiare le issues secondo criteri di interesse (come ambito o prioritá) vengono utilizzate le labels.
\\\\\textbf{Labels:}
\begin{itemize}
	\item Da approvare: indica una issue o pull request che necessita di approvazione
	\item Da verificare: indica una issue o pull request che necessita di verifica
	\item Documentazione: indica una issue o pull request riguardante la documentazione
	\item P=Bassa: indica una issue o pull request a prioritá bassa (scadenza lontana o non limita il lavoro altrui)
	\item P=Media: indica una issue o pull request a prioritá media (scadenza di almeno due settimane o non limita il lavoro altrui)
	\item P=Alta indica una issue o pull request a prioritá alta (scadenza breve o limita il lavoro altrui)
	\item Presentazione: indica una issue riguardante la redazione di una presentazione in classe o al proponente
\end{itemize}

Nel caso un membro del gruppo dovesse rendersi conto che l'issue\textsubscript{g} che sta svolgendo potrebbe essere suddiviso in ulteriori issue\textsubscript{g}, dovrà rivolgersi al responsabile, che è l'unico che può aggiungere, modificare o eliminare le issue\textsubscript{g}.\\
Per raggruppare piú issue si marca ciascuna con la label di raggruppamento.
Il nome di una label di raggruppamento viene preceduto da una G maiuscola (che sta per gruppo) seguito dal nome dell'attivitá da svolgere. esempio: "G Piano di progetto"
indica un gruppo di labels relative al documento "Piano di progetto".
Al completamento di tutte le issues di un gruppo si potrá scegliere se eliminare la label o tenerla per raggruppare issues future dello stesso gruppo.
\\\\
\textbf{Utilizzo delle checkbox:}
\\\\
Una issue può essere suddivisa in più attività tramite delle checkbox in base alla grandezza o
alla necessità di suddividere il lavoro. Al completamento di un'attività si spunta la checkbox
corrispondente in modo da avvisare il gruppo sullo stato di avanzamento di quella issue.
\subsubsection{\textit{Glossario}}
All'interno del documento si possono trovare dei termini che possono risultare ambigui a seconda del contesto, o non conosciuti dagli utilizzatori.\\\
Per ovviare ad errori di incomprensione che possono portare a problemi di vario genere e rallentamenti si è deciso di stilare un elenco di termini 
di interesse accompagnati da una descrizione dettagliata del loro significato.\\
I termini presenti all'interno del \textit{Glossario} vengono indicati con il pedice 'g' come nell'esempio seguente: termine\textsubscript{g}.
Il \textit{Glossario} ordina i termini in ordine alfabetico in modo da permetterne una facile e veloce ricerca.
Ogni componente del gruppo all'inserimento di un termine ritenuto ambiguo deve preoccuparsi di aggiornare il \textit{Glossario} in modo da mantenerlo sempre aggiornato.

\section{Organizzazione del gruppo}
\subsubsection{Ruoli}
I componenti del gruppo si suddivideranno nei seguenti ruoli per periodi di circa 2-3 settimane (dipendentemente dalle esigenze del periodo) e al termine del periodo i ruoli verranno risuddivisi. 
Visto che nelle varie fasi di sviluppo del progetto le attività da svolgere variano, non sempre sarà necessario coprire tutti i ruoli.\\
Inoltre sarà necessario tenere traccia delle ore che ogni componente dedica al progetto ed il ruolo associato a quelle ore, in modo da andare a rispettare la tabella degli impegni individuali.\\
I ruoli e le loro competenze sono i seguenti:

\paragraph{Responsabile}
Deve avere la visione d'insieme del progetto e coordinare i membri, inoltre si occupa di rappresentare il gruppo con le interazione esterne (proponente, committente ecc...). Le sue competenze specifiche sono:
\begin{itemize}
	\item Ad ogni iterazione\textsubscript{g} c'è un solo responsabile;
	\item Presentare il diario di bordo in aula;
	\item Redarre l'ordine del giorno prima di ogni meeting interno del gruppo;
	\item Suddivide le attività del gruppo in singole issue (ma non le assegna ai membri del gruppo);
	\item In fase di release\textsubscript{g} si occupa di approvare\textsubscript{g} tutti i documenti che necessitano approvazione.
\end{itemize}

\paragraph{Analista}
Si occupa di trasformare i bisogni del proponente nelle aspettative che il gruppo deve soddisfare per sviluppare un prodotto professionale. Le sue competenze specifiche sono:
\begin{itemize}
	\item Interrogare il proponente riguardo allo scopo del prodotto e le funzionalità che deve avere;
	\item Studiare le risposte del proponente per identificare i requisiti\textsubscript{g} e redarre l' \textit{Analisi dei Requisiti}.
\end{itemize}

\paragraph{Amministratore}
Si occupa del funzionamento, mantenimento e sviluppo degli strumenti tecnologici usati dal gruppo. Le sue competenze specifiche sono:
\begin{itemize}
	\item Ad ogni iterazione\textsubscript{g} basta un solo amministratore;
	\item Gestione delle segnalazioni e problemi dei membri del gruppo riguardanti problemi e malfunzionamenti con gli strumenti tecnologici;
	\item Valuta l'utilizzo di nuove tecnologie e ne fa uno studio preliminare per poter presentare al gruppo i pro e i contro del suo utilizzo.
\end{itemize}

\paragraph{Progettista}
Si occupa di scegliere la modalità migliore per soddisfare le aspettative del committente che gli analisti hanno ricavato dall'analisi dei requisiti\textsubscript{g}. Le sue competenze specifiche sono:
\begin{itemize}
	\item Scegliere eventuali pattern architetturali da implementare;
	\item Sviluppare lo schema UML\textsubscript{g} delle classi\textsubscript{g}.
\end{itemize}

\paragraph{Programmatore}
Si occupa di implementare le scelte e i modelli fatti dal progettista. Le sue competenze specifiche sono:
\begin{itemize}
	\item Scrivere il codice atto a implementare lo schema delle classi;
	\item Scrivere eventuali test;
	\item Scrivere la documentazione per la comprensione del codice che scrive.
\end{itemize}

\paragraph{Verificatore}
Si occupa di controllare che ogni file che viene caricato in un branch protetto della repository\textsubscript{g} sia conforme alle norme di progetto. Le sue competenze specifiche sono:
\begin{itemize}
	\item Controllare i file modificati o aggiunti durante una pull request tra un ramo non protetto e un ramo protetto siano conformi alle norme di progetto e cercano errori di altra natura (ortografici, sintattici, logici, build ecc...).
\end{itemize}
\pagebreak

\end{document}
